\section{The orbit obstruction for homogeneous minors}\label{mi:section}

A non-principal $2\times2$ minor with four distinct row and column indices
has four independent entries. Write
\[
 M=\begin{pmatrix}a&b\\c&d\end{pmatrix},\qquad
 S_{\det}=\{M:\det M=0\}.
\]
Every rank-one matrix is $uv^T$, so this is exactly the projection of the
outer-product equations onto those entries. After a column interchange,
we may assume $\det\bar M>0$. Throughout this section $\rc>0$, the projected
rays are matrices $P_j$, and the finite corner bound is
\[
 \zK=\min\left\{\rc^T\lambda:
       \lambda\ge0,\ \det\left(\bar M+\sum_j\lambda_jP_j\right)=0\right\}.
\]
For $z\ge0$, put
$T_z=\conv\{\bar M,\bar M+(z/\omega_j)P_j:j=1,\ldots,N\}$.
The determinant has signature $(2,2)$: the coordinates
\[
 x=\left(\frac{a+d}{2},\frac{c-b}{2}\right),\qquad
 y=\left(\frac{d-a}{2},\frac{b+c}{2}\right)
 \quad\hbox{give}\quad \det M=\norm{x}^2-\norm{y}^2.
\]
Its polar form is $B(M,N)=\tfrac12\tr(\adj(M)N)$.

\begin{lemma}\label{mi:free-cones}
Let $C$ be closed and convex with $\bar M\in\intt C$.
Then $C$ is $S_{\det}$-free if and only if $C\subseteq\{\det\ge0\}$,
and if and only if it is $\{\det\le0\}$-free. Every maximal such set is
a cone. The corner bounds for $\{\det=0\}$ and $\{\det\le0\}$ agree.
\end{lemma}
\begin{proof}
The determinant is positive on the connected set $\intt C$ if it never
vanishes there, and $C=\cl(\intt C)$. Conversely, every singular matrix
has matrices of negative determinant arbitrarily close to it: use a
direction on which its nonzero determinant gradient is negative, or
$\diag(1,-1)$ at zero. Thus $\intt\{\det\ge0\}=\{\det>0\}$.
Homogeneity gives $\cl(\cone C)\subseteq\{\det\ge0\}$; this closed convex
cone is free and contains $C$, proving the maximality assertion.
Finally a negative determinant reached at $\lambda\ge0$ must first vanish
on the segment from zero to $\lambda$, at no greater objective cost.
\end{proof}

The polar-factor choice used by the minor intersection-cut formula
\citep{BCM2020,ChmielaMunozSerrano2020ZIB} is a member of a larger finite-dimensional family.
Define
\[
 C_F=\{M:\sym(F^TM)\succeq0\},\qquad \det F>0.
\]

\begin{proposition}[Polar-factor choice]\label{mi:polar}
Let $\bar M=UH$ be the polar decomposition, where $U\in SO(2)$ and
$H\succ0$ is symmetric. The set
\[
 \{M:\norm{y(M)}\le x(\bar M)^Tx(M)/\norm{x(\bar M)}\}
\]
equals $C_U$. It is the rotation-family member chosen by
$(\lambda_1,\lambda_2)\propto(\bar a+\bar d,\bar b-\bar c)$ in the
outer-product-free inequality
$\lambda_1(a+d)+\lambda_2(b-c)\ge\norm{(b+c,a-d)}$.
For $\det\bar M<0$, a column interchange yields $C_W$, with $W$ the
orthogonal polar factor of $\bar M$ and $\det W=-1$.
\end{proposition}
\begin{proof}
For the rotation $R_\phi$, the condition
$\sym(R_\phi^TM)\succeq0$ is
$(\cos\phi,\sin\phi)^Tx(M)\ge\norm{y(M)}$.
Choose $(\cos\phi,\sin\phi)=x(\bar M)/\norm{x(\bar M)}$.
Then $R_\phi^T\bar M$ is symmetric, has positive trace and positive
determinant, and is therefore positive definite. Uniqueness of the polar
decomposition gives $R_\phi=U$. The stated outer-product-free inequality
has precisely these coordinates. If $\Pi$ interchanges columns, the polar
decomposition of $\bar M\Pi$ is
$(W\Pi)(\Pi^TH\Pi)$, and congruence by $\Pi$ gives the last assertion.
\end{proof}

\begin{proposition}[Orbit and point rule]\label{mi:orbit}
Let $G_+$ consist of $M\mapsto AMB^T$ and $M\mapsto AM^TB^T$, with
$A,B\in GL_2(\R)$ and $\det A\det B>0$.
\begin{enumerate}
\item Each $C_F$ is a maximal $S_{\det}$-free cone, and
$C_F=C_{F'}$ exactly when $F'=cF$ for $c>0$.
The $G_+$-orbit of $C_U$ is the three-parameter family of all $C_F$.
\item Under $M\mapsto AMB^T$, its parameter changes by
$F'^T=BF^TA^{-1}$. Transposition sends $C_F$ to $C_{F^{-1}}$.
\item Applying the polar-factor choice in every $G_+$ coordinate system
gives precisely
\[
 \mathcal P(\bar M)=\{C_F:F^T\bar M\text{ is symmetric positive definite}\}.
\]
This two-parameter family is equivariant under $G_+$.
\item The one-parameter rotation family $\{C_R:R\in SO(2)\}$ intersects
$\mathcal P(\bar M)$ in the single set $C_U$.
\end{enumerate}
\end{proposition}
\noindent
The proof in \cref{mi:orbit-proof} includes a direct maximality argument.
Thus these homogeneous orbit sets need no maximal-completion step.
The full family of maximal determinant-free sets can be larger than this orbit.
The general homogeneous classification uses nonexpansive sphere maps
$\Gamma:S^1\to S^1$, subject to its convex-hull condition
\citep{MPS2025}; the orbit is a finite-dimensional subfamily.

The best member of each of these families has a small semidefinite
feasibility formulation. The parameter spaces for $F^T$ are, respectively,
$\R^{2\times2}$, $\{G:G\bar M\text{ is symmetric}\}$,
$\operatorname{span}\{I,J\}$, and $\R U^T$, where
$J=\begin{psmallmatrix}0&1\\-1&0\end{psmallmatrix}$.

\begin{proposition}[Family optimization and exact certificates]\label{mi:certificates}
For any of these linear spaces $L$, a parameter $G\in L$ gives a family
member with $z_{C_F}(\rc)\ge z$ if and only if
\[
 \sym(G\bar M)\succ0,\qquad \sym(GV)\succeq0\quad(V\text{ a vertex of }T_z),
 \qquad G=F^T.
\]
Feasibility is monotone in $z$ and uses at most $N+1$ symmetric
$2\times2$ blocks and four scalar variables.
An upper certificate consists of positive definite symmetric matrices
$Y_V$, one for each vertex including $\bar M$, such that
\begin{equation}\label{mi:dual-certificate}
 \sum_V\tr(GVY_V)=0\qquad\text{for every }G\in L.
\end{equation}
Its existence implies $\sup_{F^T\in L}z_{C_F}(\rc)\le z$.
\end{proposition}
\begin{proof}
The first conditions are exactly $\bar M\in\intt C_F$ and $T_z\subseteq C_F$;
positive definiteness at the apex also forces $\det F>0$.
Smaller simplices are contained in $T_z$. If a feasible $G$ and the
certificate both existed, the left side of \cref{mi:dual-certificate}
would be a sum of nonnegative matrix inner products, with the apex term
strictly positive. This contradicts its being zero.
\end{proof}

\begin{proposition}[Support and tangent edges]\label{mi:support}
There is a corner minimizer whose supported projected rays are linearly
independent. Every such minimizer has support at most three; support three
requires $\bar M$ to lie in the span of its three rays.
A support-two minimizer lies in the relative interior of an edge of
$T_{\zK}$, with edge direction $D_{\mathrm{edge}}$ satisfying
$\tr(\adj(M_*)D_{\mathrm{edge}})=0$ and $\det D_{\mathrm{edge}}\ge0$.
\end{proposition}
\noindent
The proof in \cref{mi:support-proof} uses the signature $(2,2)$ and keeps
the independence hypothesis explicit. Duplicating one supported ray can
produce other minimizers with larger support.

A tangent edge with $\det D_{\mathrm{edge}}>0$ constrains the entire orbit parameter to
one projective degree of freedom.

\begin{lemma}[Tangent-edge pencil]\label{mi:pencil}
Let $M_*=ab^T\ne0$ and let $D_{\mathrm{edge}}$ satisfy
$\tr(\adj(M_*)D_{\mathrm{edge}})=0$, $\det D_{\mathrm{edge}}>0$. Set
$G_0=J\adj(D_{\mathrm{edge}})$ and $G_1=J\adj(M_*)$.
Then $G_1a=0$, $G_0a=c_0b$ for some $c_0\ne0$, and
$\det(G_0+tG_1)=\det D_{\mathrm{edge}}$ for all $t$.
If an orbit set contains $M_*+sD_{\mathrm{edge}}$ for all $|s|\le\varepsilon$, then
\[
 F^T=\theta(G_0+tG_1),\qquad \theta c_0>0.
\]
\end{lemma}
\noindent
Consequently each additional vertex imposes an interval condition on $t$.
Disjoint intervals obstruct attainment, but do not by themselves prove a
strict supremum gap. The following rational certificates do prove one.

\begin{theorem}[Full-dimensional minor counterexamples]\label{mi:examples}
The rational four-ray corners A, B and S1 in \cref{mi:example-data} have
$\rc=(1,1,1,1)$, independent rays, and $\zK=1$ with a unique minimizer.
A and B have tangent-edge minimizers; S1 has a support-one minimizer with
strictly positive inactive-ray multipliers and transverse incident edges.
Their orbit suprema satisfy
\[
 z_{\mathrm{orbit}}(\mathrm A)\le\frac9{20},\qquad
 z_{\mathrm{orbit}}(\mathrm B)\le\frac{461}{500},\qquad
 z_{\mathrm{orbit}}(\mathrm{S1})\le\frac{779}{1000}.
\]
For A, an unrestricted maximal free set does attain the corner bound one.
The strict gap for A persists on an open neighborhood of its apex, rays
and positive objective weights.
\end{theorem}
\noindent
\Cref{mi:example-proof} gives exact face calculations and the complete
rational dual matrices. In particular, support one and transversality
alone do not guarantee orbit exactness.

\begin{corollary}[Bilinear slices and exact comparisons]\label{mi:embedding}
Embed a bilinear corner by
$M=\begin{psmallmatrix}w&x\\y&h\end{psmallmatrix}$,
with apex $(\bar w,\bar x,\bar y,1)$ and rays
$(p_{jw},p_{jx},p_{jy},0)$.
Its corner bound and minor orbit bound equal, respectively, the bilinear
corner bound and family-$\famA$ bound.
For the three rational bilinear examples of
\cref{fd:counterexample,fd:support-one}, including the second tangent-edge
instance, exact orbit brackets are
$[0.9753853,0.9753864]$, $[0.8347663,0.8347772]$, and
$[0.9838465,0.9838468]$.
Thus A and S1 have larger gaps than these three selected bilinear examples,
whereas B has a smaller gap than the second bilinear example.
\end{corollary}
\begin{proof}
On $h=1$, $\det M=w-xy$, so \cref{mi:free-cones} identifies the corner
bounds. The slice of $C_F$ has exactly the bilinear family-$\famA$ ray steps.
The rational parameters and dual matrices in
\cref{mi:comparison-certificates} prove both ends of the displayed brackets
and all the strict comparisons.
\end{proof}

There is therefore no ordering of the gaps across all minor and bilinear
corners. Since every bilinear corner embeds as a minor corner, the infimum
minor ratio is no greater than the infimum bilinear ratio; this does not
compare full-dimensional corners under specified conditioning bounds.

\begin{proposition}[Contact restriction]\label{mi:contact}
If a closed convex free set $C$ with $\bar M\in\intt C$ satisfies
$z_C=\zK<\infty$, then $T_{\zK}\subseteq C$ and every corner minimizer
belongs to $C\cap S_{\det}$. For an orbit set,
\[
 C_F\cap S_{\det}=\{F^{-T}bb^T:b\in\R^2\}.
\]
This contact cone has dimension two, whereas $S_{\det}$ has dimension three.
If the unique minimizer is on ray $j$, the polar-factor set can be exact
only if $x(P_j)$ is parallel to $x(\bar M)$, a single linear condition
on that ray.
\end{proposition}
\begin{proof}
Exactness implies $\alpha_j\ge\zK/\omega_j$ for every ray, hence simplex
containment. For a rank-one $ab^T$, positive semidefiniteness of
$\sym((F^Ta)b^T)$ is equivalent to $F^Ta=cb$ with $c\ge0$.
Absorb $c$ into $b$ to obtain the contact formula.
At a support-one contact, $U^TM_*$ and $U^T\bar M$ are symmetric;
their difference forces the skew coordinate of $U^TP_j$ to vanish,
which is the stated condition.
\end{proof}

\begin{proposition}[Variable scaling]\label{mi:scaling}
For $0<\delta\le1/4$, take
\[
 \bar M=\diag(1,\delta),\quad
 (P_1,P_2,P_3,P_4)=(E_{12},-E_{11},-\delta E_{22},-\delta E_{21}),
 \quad\rc=(1,1,1,1).
\]
Then $\zK=1$, the polar-factor set gives $2\sqrt\delta$, and both the
point-rule family and the orbit attain one. Every rotation-family set gives
at most $((16+2\sqrt{73})/3)\sqrt\delta$.
This is the image under $M\mapsto\diag(1,\delta)M$ of a corner where the
polar-factor choice is exact.
\end{proposition}
\noindent
The calculations in \cref{mi:scaling-proof} show that both the polar-factor
choice and the full rotation family can lose arbitrarily much under one
variable scaling, while the equivariant point-rule and orbit families remain
exact.

\begin{proposition}[Positive-determinant principal minors]\label{mi:principal}
On the symmetric matrix space, let $\det\bar M>0$.
If $\tr\bar M>0$, the positive semidefinite cone is the unique maximal
determinant-free set containing $\bar M$ in its interior; if
$\tr\bar M<0$, it is the negative semidefinite cone.
The polar-factor set restricted to this space is that cone and its
intersection cut attains $\zK$.
\end{proposition}
\begin{proof}
The two components of $\{\det>0\}$ on symmetric matrices are the positive
and negative definite cones. A connected free interior containing
$\bar M$ stays in its component, so its closure lies in the corresponding
semidefinite cone. That cone is convex and free, proving uniqueness.
Every $T_z$ with $z<\zK$ stays in the same component, so the cone's ray
steps give $z_C\ge\zK$; validity gives the reverse inequality.
The polar factors are respectively $I$ and $-I$.
\end{proof}

These conclusions use the full rank-one equation. If one entry is diagonal,
the outer-product projection instead has closure
$\{\det M=0,a\ge0\}$, which can permit larger free sets.
Negative-determinant principal minors and such bound-aware enlargements
therefore require a separate analysis; the precise projection and the
principal-minor distinction are recorded in \cref{mi:other-minors}.
